% Part: first-order-logic
% Chapter: introduction
% Section: semantic-notions

\documentclass[../../../include/open-logic-section]{subfiles}

\begin{document}

\olfileid{fol}{int}{sem}

% SEGMENT OLP-0145-S01
\olsection{பொருண்மையியல் கருத்துகள்}

முன்னர் $\Sat{M}{!A}[s]$ பொருந்துகிறதா என்று கருதும்போது, வசதிக்காக $s$
எல்லா !!{variable}s க்கும் மதிப்புகளை ஒதுக்குமாறு வைத்தோம்; ஆனால்~$!A$ இல்
வரும் !!{variable}s க்கு அது ஒதுக்கும் மதிப்புகள் மட்டுமே பயன்படுகின்றன.
உண்மையில்,~$!A$ இல் உள்ள \emph{கட்டற்ற} மாறிகளின் மதிப்புகளே முக்கியம்.
நாம் கவனமாக இருப்பதால், இந்த உண்மையை நிறுவப் போகிறோம். !!{sentence}s க்கு
கட்டற்ற மாறிகள் எதுவும் இல்லாததால், அவை !!a{structure} இல்
நிறைவுறுத்தப்படுகிறதா இல்லையா என்பதில்~$s$ சிறிதும் முக்கியமல்ல. ஆகவே, $!A$
!!a{sentence} ஆக இருக்கும்போது, $\Sat{M}{!A}$ என்பதை
``எல்லா~$s$ களுக்கும் $\Sat{M}{!A}[s]$'' என்று வரையறுக்கலாம்; நிகழ்வாக,
இது குறைந்தது ஒரு~$s$ க்கு $\Sat{M}{!A}[s]$ என்பதற்கு அப்போதும் அப்போது
மட்டுமே பொருந்தும். !!{formula}s க்கான நிறைவுறுத்தலைச் செயல்படத்தக்கவாறு
வரையறுக்க !!{variable} ஒதுக்கீடுகளை அறிமுகப்படுத்த வேண்டும்; ஆனால்
!!{sentence}s க்கு நிறைவுறுத்தல் !!{variable} ஒதுக்கீடுகளைச் சார்ந்ததல்ல.

ஒருமுறை ``$\Sat{M}{!A}$'' என்பதன் வரையறை கிடைத்ததும், முதல்தர தருக்கத்தின்
!!{sentence}s குறித்து ``நிலை'' மற்றும் ``இதில் மெய்'' என்பவற்றின் பொருள்
நமக்குத் தெரியும். $\Sat{M}{!A}$ இன் வரையறையை அடிப்படையாகக் கொண்டு,
செல்லுபடியாகும் தன்மை, உட்கிடை, நிறைவுசெய்யத்தக்க தன்மை ஆகிய அடிப்படைப்
பொருண்மையியல் கருத்துகளை வரையறுக்கலாம். ஒரு !!a{sentence} செல்லுபடியானது,
$\Entails !A$, எனில் ஒவ்வொரு !!{structure} உம் அதை நிறைவுறுத்தும்.
!!{sentence}s கொண்ட ஒரு கணத்திலிருந்து அது உட்கிடையாகிறது,
$\Gamma \Entails !A$, எனில்~$\Gamma$ இல் உள்ள எல்லா !!{sentence}s ஐயும்
நிறைவுறுத்தும் ஒவ்வொரு !!{structure} உம்~$!A$ ஐயும் நிறைவுறுத்தும். மேலும்,
!!{sentence}s கொண்ட கணம் ஒன்று நிறைவுசெய்யத்தக்கது எனில், அதிலுள்ள எல்லா
!!{sentence}s ஐயும் ஒரே நேரத்தில் நிறைவுறுத்தும் !!a{structure} ஒன்று உள்ளது.

!!{formula}s தொகுத்தறிதலால் வரையறுக்கப்பட்டுள்ளன; நிறைவுறுத்தலும்
!!{formula}s இன் கட்டமைப்பின் மீது தொகுத்தறிதலால் வரையறுக்கப்படுகிறது.
எனவே நமது பொருண்மையியலின் பண்புகளை நிறுவவும், வரையறுக்கப்பட்ட
பொருண்மையியல் கருத்துகளை ஒன்றோடொன்று தொடர்புபடுத்தவும் தொகுத்தறிதலைப்
பயன்படுத்தலாம். சில பண்புகளைத் தனித்தனியாக ஆர்வமூட்டுவதால் மட்டும் அல்லாமல்,
பொருண்மையியலுக்கும் !!{derivation} முறைமைகளுக்கும் இடையிலான உறவை ஆராயும்
போது அவற்றில் பல உதவியாக இருக்கும் என்பதாலும் அவற்றைச் சேகரித்து நிறுவுவோம்.
இதற்காக, இதுவரை குறிப்பிடாத சில தொடரமைப்புக் கருத்துகளையும் செயல்களையும்
(துல்லியமாக, தொகுத்தறிதலால்) வரையறுக்க வேண்டும்.

\end{document}
